% =============================================================================
%                    WEEK 4: NEURAL LMs, RNNs, AND LSTMs
% =============================================================================
\section{Week 4: Neural Language Models, RNNs, and LSTMs}

% -----------------------------------------------------------------------------
%                              4.1 TOPIC SUMMARY
% -----------------------------------------------------------------------------
\subsection{Topic Summary}

\begin{keybox}[Learning Objectives]
By the end of this week, you will be able to:
\begin{enumerate}
  \item Explain \textbf{embeddings} and the \textbf{embedding matrix}, including how token indices (or one-hot vectors) are mapped to embedding vectors.
  \item Use embeddings as input to \textbf{neural classifiers}, and combine them with methods such as \textbf{pooling} to produce sentence-level representations.
  \item Describe simple \textbf{feedforward neural language models}, which use token embeddings in a fixed window to predict the next word.
  \item Compare neural LMs with n-gram models, highlighting advantages such as \textbf{generalisation} and longer context.
  \item Explain \textbf{recurrent neural networks (RNNs)} and how they process sequences with hidden states.
  \item Describe how RNNs are trained, including \textbf{backpropagation through time} and \textbf{teacher forcing}.
  \item Apply RNNs to language modelling and NLP tasks, such as sequence labelling and text generation.
  \item Motivate and explain \textbf{LSTMs}, including their \textbf{gating mechanisms} for handling long-distance dependencies.
\end{enumerate}
\end{keybox}

\separator

\subsubsection*{The Big Picture}

This week transitions from count-based and simple neural methods to \textbf{sequence-aware architectures}. We start with embeddings as learned representations fed into feedforward classifiers and language models, then introduce \textbf{RNNs} which maintain hidden states across time steps, enabling modelling of sequences of arbitrary length. Finally, we cover \textbf{LSTMs}, which solve the vanishing gradient problem through gating mechanisms, enabling capture of long-distance dependencies.

\separator

\subsubsection*{What Examiners Like to Ask}

\begin{itemize}
  \item Explain how the \textbf{embedding matrix} $E$ maps one-hot vectors (or indices) to dense embeddings.
  \item Describe \textbf{pooling} (mean pooling) for converting variable-length sequences into fixed-size representations.
  \item Compare \textbf{feedforward neural LMs} to \textbf{n-gram LMs}: context size, generalisation, sparsity handling.
  \item Write the \textbf{RNN forward equations} ($h_t = g(Uh_{t-1} + Wx_t)$, $y_t = f(Vh_t)$).
  \item Explain \textbf{unrolling} for training and why RNNs require incremental inference.
  \item Define \textbf{teacher forcing} and \textbf{weight tying} in RNN language model training.
  \item Describe the \textbf{vanishing gradient problem} and why it motivates LSTMs.
  \item Explain the \textbf{LSTM gating mechanisms}: forget gate, input/add gate, output gate.
  \item Describe the \textbf{encoder-decoder architecture} and its components.
  \item Explain \textbf{attention} as a solution to the context bottleneck in encoder-decoder models.
\end{itemize}

\separator

% -----------------------------------------------------------------------------
%                 4.2 OVERVIEW + CORE CONCEPTS + EXTENDED THEORY
% -----------------------------------------------------------------------------
\subsection{Overview + Core Concepts + Extended Theory}

\subsubsection*{Roadmap}
\begin{enumerate}
  \item Embeddings and the embedding matrix $E$
  \item Embeddings as input to feedforward classifiers (pooling)
  \item Feedforward neural language models
  \item Neural LMs vs.\ n-gram models
  \item Recurrent Neural Networks (RNNs / Elman Nets)
  \item Training RNNs: unrolling, backpropagation through time, teacher forcing
  \item RNNs for NLP tasks: LM, sequence labelling, classification, generation
  \item LSTMs: gating mechanisms for long-distance dependencies
  \item Encoder-decoder architecture
  \item Attention mechanism
\end{enumerate}

\bigseparator

% --- 4.2.1 Embeddings and the Embedding Matrix ---
\subsubsection{Embeddings and the Embedding Matrix}

\begin{defbox}[Embedding Matrix $E$]
An \textbf{embedding} is a dense vector of dimension $d$ representing a token.

The \textbf{embedding matrix} $E$ is a learned lookup table:
\begin{itemize}
  \item Shape: $|V| \times d$ (one row per vocabulary token)
  \item Each row $E[i]$ is the embedding for token $i$
\end{itemize}

Embedding matrices are \textbf{central to all modern NLP systems}.
\end{defbox}

\begin{defbox}[One-Hot to Embedding Conversion]
\textbf{Method 1: Indexing} (efficient)
\begin{itemize}
  \item Given token index $i$, select row $E[i]$ directly.
\end{itemize}

\textbf{Method 2: Matrix multiplication} (conceptual)
\begin{itemize}
  \item Represent token as one-hot vector $\vec{x} \in \mathbb{R}^{|V|}$ where $x_i = 1$, all others $0$.
  \item Multiply: $\vec{e} = \vec{x}^\top E$ selects row $i$.
\end{itemize}

Both are equivalent; indexing is faster in practice.
\end{defbox}

\begin{exbox}[Example: Selecting Embeddings]
\textbf{Vocabulary:} $|V| = 10000$, embedding dimension $d = 300$

\textbf{Token:} ``dessert'' has index $i = 3$

\textbf{One-hot vector:} $\vec{x} = [0, 0, 1, 0, \ldots, 0]$ (1 at position 3)

\textbf{Embedding:} $\vec{e}_{\text{dessert}} = E[3] \in \mathbb{R}^{300}$

For $N$ tokens, we get $N$ embeddings, forming an $N \times d$ matrix.
\end{exbox}

\bigseparator

% --- 4.2.2 Embeddings as Input to Classifiers ---
\subsubsection{Embeddings as Input to Feedforward Classifiers}

\begin{keybox}[Representation Learning]
Instead of hand-crafted features (like ``count of positive lexicon words''), we use \textbf{learned embeddings} as input features.

\textbf{Benefit:} The network learns useful representations from data.
\end{keybox}

\textbf{Problem:} Text comes in variable lengths, but classifiers need fixed-size input.

\textbf{Solutions:}
\begin{center}
\renewcommand{\arraystretch}{1.3}
\begin{tabular}{@{} l p{9cm} @{}}
\toprule
\textbf{Method} & \textbf{Description} \\
\midrule
Concatenation & Concatenate all $N$ embeddings into one vector of shape $[1 \times Nd]$. Pad shorter sequences; truncate longer ones. \\
Pooling & Combine $N$ embeddings into a single $[1 \times d]$ vector (same dimensionality as one word). \\
\bottomrule
\end{tabular}
\end{center}

\begin{defbox}[Mean Pooling]
Given $N$ token embeddings $e(w_1), \ldots, e(w_N)$, compute the element-wise mean:
\[
\vec{x}_{\text{mean}} = \frac{1}{N} \sum_{i=1}^{N} e(w_i)
\]

\textbf{Intuition:} For sentiment analysis, exact word position matters less than overall word content.

\textbf{Equations for pooled classifier:}
\begin{align*}
\vec{x} &= \text{mean}(e(w_1), \ldots, e(w_N)) \\
\vec{h} &= \sigma(W\vec{x} + \vec{b}) \\
\vec{z} &= U\vec{h} \\
\hat{\vec{y}} &= \text{softmax}(\vec{z})
\end{align*}
\end{defbox}

\begin{tipbox}[Pooling vs.\ Concatenation Trade-off]
\begin{itemize}
  \item \textbf{Pooling:} Less information (loses position), but efficient and handles variable length naturally.
  \item \textbf{Concatenation:} Preserves position, but requires padding/truncation and more parameters.
\end{itemize}
\end{tipbox}

\bigseparator

% --- 4.2.3 Feedforward Neural Language Models ---
\subsubsection{Feedforward Neural Language Models}

\begin{defbox}[Feedforward Neural LM]
\textbf{Task:} Predict next word $w_t$ given prior words $w_{t-1}, w_{t-2}, \ldots$

\textbf{Approach:} Use a \textbf{sliding window} of fixed size $N-1$ (like n-grams).

\textbf{Architecture:}
\begin{enumerate}
  \item Look up embeddings for context words: $e_{t-1}, e_{t-2}, \ldots, e_{t-N+1}$
  \item Concatenate into input vector: $\vec{e} = [e_{t-N+1}; \ldots; e_{t-1}]$
  \item Pass through hidden layer: $\vec{h} = \sigma(W\vec{e} + \vec{b})$
  \item Output layer with softmax over vocabulary: $\hat{\vec{y}} = \text{softmax}(U\vec{h})$
\end{enumerate}
\end{defbox}

\begin{exbox}[Example: 4-gram Feedforward LM]
\textbf{Context:} ``and thanks for all the \underline{~~~}''

\textbf{Window ($N=4$):} [``for'', ``all'', ``the''] $\to$ predict ``fish''

\textbf{Dimensions:}
\begin{itemize}
  \item Embedding layer: $[1 \times Nd]$ (concatenated embeddings)
  \item Hidden layer: $[1 \times d_h]$
  \item Output layer: $[1 \times |V|]$ (probability over vocabulary)
\end{itemize}

$\hat{y}_{42}$ gives $P(w_t = \text{``fish''} | w_{t-3}, w_{t-2}, w_{t-1})$
\end{exbox}

\bigseparator

% --- 4.2.4 Neural LMs vs N-gram LMs ---
\subsubsection{Neural LMs vs.\ N-gram LMs}

\begin{thmbox}[Why Neural LMs Outperform N-gram LMs]
\textbf{Generalisation through embeddings:}

\textbf{Training data:} ``I have to make sure that the \textbf{cat} gets fed.''

\textbf{Test data:} ``I forgot to make sure that the \textbf{dog} gets \underline{~~~}''

\begin{itemize}
  \item \textbf{N-gram LM:} Never saw ``dog gets''; cannot predict ``fed''.
  \item \textbf{Neural LM:} Uses similarity of ``cat'' and ``dog'' embeddings to generalise and predict ``fed''.
\end{itemize}
\end{thmbox}

\begin{center}
\renewcommand{\arraystretch}{1.3}
\begin{tabular}{@{} l p{5cm} p{5cm} @{}}
\toprule
& \textbf{N-gram LM} & \textbf{Neural LM} \\
\midrule
Representation & Discrete word identities & Continuous embeddings \\
Generalisation & None (exact match only) & Similar words $\to$ similar predictions \\
Sparsity & Severe (many zero counts) & No zeros (dense vectors) \\
Context & Fixed $N-1$ words & Fixed window (feedforward) \\
Parameters & $O(|V|^N)$ & $O(|V| \cdot d + d \cdot d_h)$ \\
\bottomrule
\end{tabular}
\end{center}

\bigseparator

% --- 4.2.5 Recurrent Neural Networks (RNNs) ---
\subsubsection{Recurrent Neural Networks (RNNs)}

\begin{defbox}[Recurrent Neural Network (RNN)]
A \textbf{recurrent neural network} is any network containing a \textbf{cycle} in its connections.

The hidden state $h_t$ depends on:
\begin{itemize}
  \item The current input $x_t$
  \item The previous hidden state $h_{t-1}$
\end{itemize}

This allows information to \textbf{persist across time steps}.
\end{defbox}

\begin{defbox}[Simple RNN (Elman Net) Equations]
\textbf{Forward inference:}
\begin{align*}
h_t &= g(U h_{t-1} + W x_t) \\
y_t &= f(V h_t)
\end{align*}

where:
\begin{itemize}
  \item $x_t$: input at time $t$ (e.g., embedding)
  \item $h_t$: hidden state at time $t$
  \item $y_t$: output at time $t$
  \item $W$: input-to-hidden weights $[d_h \times d]$
  \item $U$: hidden-to-hidden weights $[d_h \times d_h]$
  \item $V$: hidden-to-output weights $[|V| \times d_h]$
  \item $g$: activation function (e.g., tanh, ReLU)
  \item $f$: output function (e.g., softmax)
\end{itemize}
\end{defbox}

\begin{tipbox}[Key Property: Sequential Inference]
Computing $h_t$ \textbf{requires} $h_{t-1}$.

$\Rightarrow$ Inference must proceed \textbf{incrementally} from start to end.

$\Rightarrow$ Cannot parallelise across time steps (unlike feedforward networks).
\end{tipbox}

\begin{exbox}[RNN Forward Inference Algorithm]
\begin{verbatim}
function FORWARD_RNN(x, network):
    h_0 = 0  # initial hidden state
    for i = 1 to LENGTH(x):
        h_i = g(U * h_{i-1} + W * x_i)
        y_i = f(V * h_i)
    return y
\end{verbatim}

\textbf{Note:} Matrices $U$, $V$, $W$ are \textbf{shared across all time steps}.
\end{exbox}

\bigseparator

% --- 4.2.6 RNNs as Language Models ---
\subsubsection{RNNs as Language Models}

\begin{keybox}[RNN LM vs.\ Feedforward LM]
\begin{center}
\renewcommand{\arraystretch}{1.3}
\begin{tabular}{@{} l p{5cm} p{5cm} @{}}
\toprule
& \textbf{Feedforward LM} & \textbf{RNN LM} \\
\midrule
Context & Fixed window size & Entire preceding history (via $h_{t-1}$) \\
Representation & Concatenated embeddings & Hidden state encodes history \\
\bottomrule
\end{tabular}
\end{center}
\end{keybox}

\begin{defbox}[RNN Language Model Equations]
\begin{align*}
e_t &= E x_t & \text{(embedding lookup)} \\
h_t &= g(U h_{t-1} + W e_t) & \text{(update hidden state)} \\
\hat{y}_t &= \text{softmax}(V h_t) & \text{(predict next word)}
\end{align*}

\textbf{Probability of next word being $k$:}
\[
P(w_{t+1} = k \mid w_1, \ldots, w_t) = \hat{y}_t[k]
\]

\textbf{Probability of entire sequence:}
\[
P(w_{1:n}) = \prod_{i=1}^{n} P(w_i \mid w_{<i}) = \prod_{i=1}^{n} \hat{y}_{i-1}[w_i]
\]
\end{defbox}

\bigseparator

% --- 4.2.7 Training RNNs ---
\subsubsection{Training RNNs}

\begin{defbox}[Unrolling Through Time]
To train an RNN, we \textbf{unroll} it into a feedforward computational graph:
\begin{enumerate}
  \item Given input sequence, create copies of the network for each time step.
  \item Weight matrices ($U$, $V$, $W$) are \textbf{shared} across all copies.
  \item Apply standard backpropagation to the unrolled graph.
\end{enumerate}

This is called \textbf{Backpropagation Through Time (BPTT)}.
\end{defbox}

\begin{defbox}[Self-Supervision for LM Training]
\textbf{Self-supervision:} The text itself provides the training signal.
\begin{itemize}
  \item At each time step $t$, the model predicts $w_{t+1}$.
  \item The true next word $w_{t+1}$ serves as the gold label.
  \item No manual annotation required.
\end{itemize}

\textbf{Loss function:} Cross-entropy at each position.
\[
L_{CE}(\hat{y}_t, y_t) = -\log \hat{y}_t[w_{t+1}]
\]

Since $y_t$ is one-hot, only the probability of the correct word matters.
\end{defbox}

\begin{defbox}[Teacher Forcing]
\textbf{Teacher forcing:} During training, always feed the \textbf{correct (gold) previous word} as input, not the model's prediction.

\textbf{Why?}
\begin{itemize}
  \item Early in training, predictions are poor.
  \item Feeding wrong predictions would compound errors.
  \item Gold context ensures stable training signal.
\end{itemize}

\textbf{At test time:} We must use the model's own predictions (no gold available).
\end{defbox}

\begin{defbox}[Weight Tying]
The embedding matrix $E$ and output matrix $V$ serve similar roles:
\begin{itemize}
  \item $E$: maps words to embeddings (rows are word vectors)
  \item $V$: maps hidden states to word scores (rows are word vectors)
\end{itemize}

\textbf{Weight tying:} Use $V = E^\top$ instead of learning separate matrices.

\textbf{Weight-tied RNN LM equations:}
\begin{align*}
e_t &= E x_t \\
h_t &= g(U h_{t-1} + W e_t) \\
\hat{y}_t &= \text{softmax}(E^\top h_t)
\end{align*}

\textbf{Benefits:} Fewer parameters, often better generalisation.
\end{defbox}

\bigseparator

% --- 4.2.8 RNNs for NLP Tasks ---
\subsubsection{RNNs for NLP Tasks}

\textbf{Four main architectures:}

\begin{center}
\renewcommand{\arraystretch}{1.3}
\begin{tabular}{@{} l p{4cm} p{6cm} @{}}
\toprule
\textbf{Task Type} & \textbf{Example} & \textbf{Architecture} \\
\midrule
Sequence Labelling & POS tagging, NER & One output per input token \\
Sequence Classification & Sentiment analysis & One output for entire sequence (use final $h_n$ or pooling) \\
Language Modelling & Next word prediction & Output at each step predicts next token \\
Encoder-Decoder & Translation, Summarisation & Separate encoder and decoder RNNs \\
\bottomrule
\end{tabular}
\end{center}

\begin{defbox}[Stacked RNNs]
Multiple RNN layers stacked on top of each other:
\begin{itemize}
  \item Output of layer 1 becomes input to layer 2.
  \item Captures more abstract representations at higher layers.
\end{itemize}
\end{defbox}

\begin{defbox}[Bidirectional RNNs]
Run \textbf{two separate RNNs}:
\begin{itemize}
  \item \textbf{Forward RNN:} processes $x_1, \ldots, x_n$ (left-to-right)
  \item \textbf{Backward RNN:} processes $x_n, \ldots, x_1$ (right-to-left)
\end{itemize}

\textbf{Concatenate} their hidden states:
\[
h_t = [\overrightarrow{h}_t ; \overleftarrow{h}_t]
\]

\textbf{Benefit:} Each position has context from \textbf{both} left and right.

\textbf{Use case:} Sequence labelling, classification (but \textbf{not} language modelling, which must be left-to-right).
\end{defbox}

\begin{exbox}[Autoregressive Generation]
To generate text:
\begin{enumerate}
  \item Start with \code{<s>} token.
  \item Predict distribution over next word.
  \item \textbf{Sample} from distribution (or take argmax for greedy decoding).
  \item Feed sampled word as next input.
  \item Repeat until \code{</s>} or max length.
\end{enumerate}
\end{exbox}

\bigseparator

% --- 4.2.9 LSTMs ---
\subsubsection{Long Short-Term Memory Networks (LSTMs)}

\begin{tipbox}[Motivation: Problems with Vanilla RNNs]
\textbf{Problem 1: Long-distance dependencies}
\begin{itemize}
  \item Example: ``The flights the airline was canceling \textbf{were} full.''
  \item ``were'' depends on ``flights'' (far away), not ``canceling'' (nearby).
  \item Hidden state must carry information over many steps.
\end{itemize}

\textbf{Problem 2: Vanishing gradients}
\begin{itemize}
  \item During backpropagation, gradients are multiplied through many time steps.
  \item Gradients can shrink exponentially, preventing learning of long-range dependencies.
\end{itemize}
\end{tipbox}

\begin{defbox}[LSTM Key Idea]
LSTMs divide context management into two subproblems:
\begin{enumerate}
  \item \textbf{Removing} information no longer needed
  \item \textbf{Adding} information likely needed for future decisions
\end{enumerate}

LSTMs introduce:
\begin{itemize}
  \item An explicit \textbf{context/cell state} $c_t$ (separate from hidden state $h_t$)
  \item \textbf{Gates}: neural circuits that control information flow
\end{itemize}
\end{defbox}

\begin{defbox}[LSTM Gates]
\textbf{1. Forget Gate} --- decides what to \textbf{remove} from context:
\begin{align*}
f_t &= \sigma(U_f h_{t-1} + W_f x_t) \\
k_t &= c_{t-1} \odot f_t
\end{align*}

\textbf{2. Input/Add Gate} --- decides what to \textbf{add} to context:
\begin{align*}
g_t &= \tanh(U_g h_{t-1} + W_g x_t) & \text{(candidate values)} \\
i_t &= \sigma(U_i h_{t-1} + W_i x_t) & \text{(gate)} \\
j_t &= g_t \odot i_t \\
c_t &= j_t + k_t & \text{(new context)}
\end{align*}

\textbf{3. Output Gate} --- decides what to \textbf{output} as hidden state:
\begin{align*}
o_t &= \sigma(U_o h_{t-1} + W_o x_t) \\
h_t &= o_t \odot \tanh(c_t)
\end{align*}

where $\odot$ denotes element-wise (Hadamard) multiplication.
\end{defbox}

\begin{tipbox}[Intuition for LSTM Gates]
\begin{itemize}
  \item \textbf{Forget gate $f_t$:} Values near 0 erase; values near 1 preserve.
  \item \textbf{Input gate $i_t$:} Values near 0 block new info; values near 1 allow it.
  \item \textbf{Output gate $o_t$:} Controls how much of cell state reaches hidden state.
\end{itemize}

The cell state $c_t$ acts as a \textbf{conveyor belt} carrying information across time with minimal modification (additive updates, not multiplicative).
\end{tipbox}

\bigseparator

% --- 4.2.10 Encoder-Decoder Architecture ---
\subsubsection{Encoder-Decoder Architecture}

\begin{defbox}[Encoder-Decoder Model]
For sequence-to-sequence tasks (translation, summarisation):

\textbf{Three components:}
\begin{enumerate}
  \item \textbf{Encoder:} Reads input sequence $x_{1:n}$, produces hidden states $h^e_{1:n}$.
  \item \textbf{Context vector:} $c = h^e_n$ (final encoder hidden state).
  \item \textbf{Decoder:} Takes $c$ as initial state, generates output sequence $y_{1:m}$.
\end{enumerate}
\end{defbox}

\begin{defbox}[Encoder-Decoder Equations]
\textbf{Encoder:}
\[
h^e_t = g(h^e_{t-1}, x_t)
\]

\textbf{Context:}
\[
c = h^e_n
\]

\textbf{Decoder:}
\begin{align*}
h^d_0 &= c \\
h^d_t &= g(\hat{y}_{t-1}, h^d_{t-1}, c) \\
\hat{y}_t &= \text{softmax}(h^d_t)
\end{align*}

\textbf{Conditional probability:}
\[
P(y \mid x) = \prod_{t=1}^{m} P(y_t \mid y_{1:t-1}, x)
\]
\end{defbox}

\begin{tipbox}[Bottleneck Problem]
The entire source sequence must pass through \textbf{one fixed-size vector $c$}.

For long sequences, this is a severe \textbf{information bottleneck}.

\textbf{Solution:} Attention mechanism.
\end{tipbox}

\bigseparator

% --- 4.2.11 Attention Mechanism ---
\subsubsection{Attention Mechanism}

\begin{defbox}[Attention: Key Idea]
Instead of using only the final encoder state:
\begin{itemize}
  \item Compute a \textbf{weighted average} of \textbf{all} encoder hidden states.
  \item Weights depend on \textbf{relevance} to the current decoder state.
  \item Different decoder steps attend to different parts of the input.
\end{itemize}
\end{defbox}

\begin{defbox}[Attention Computation]
At decoder time step $i$:

\textbf{1. Compute attention scores} (dot-product attention):
\[
\text{score}(h^d_{i-1}, h^e_j) = h^d_{i-1} \cdot h^e_j
\]

\textbf{2. Normalise with softmax to get weights:}
\[
\alpha_{ij} = \frac{\exp(\text{score}(h^d_{i-1}, h^e_j))}{\sum_k \exp(\text{score}(h^d_{i-1}, h^e_k))}
\]

\textbf{3. Compute context vector as weighted sum:}
\[
c_i = \sum_j \alpha_{ij} h^e_j
\]

\textbf{4. Use $c_i$ in decoder:}
\[
h^d_i = g(\hat{y}_{i-1}, h^d_{i-1}, c_i)
\]
\end{defbox}

\begin{tipbox}[Attention Interpretation]
\begin{itemize}
  \item $\alpha_{ij}$: how much decoder step $i$ ``attends to'' encoder position $j$.
  \item High $\alpha_{ij}$: encoder position $j$ is relevant for predicting output $i$.
  \item Attention weights are often visualised as alignment matrices.
\end{itemize}
\end{tipbox}

\begin{exbox}[Example: Translation with Attention]
\textbf{Source (English):} ``The green witch arrived''

\textbf{Target (Spanish):} ``Lleg\'{o} la bruja verde''

When generating ``bruja'' (witch):
\begin{itemize}
  \item High attention weight on ``witch''
  \item Lower weights on ``The'', ``green'', ``arrived''
\end{itemize}

The model learns soft word alignments automatically!
\end{exbox}

\bigseparator

% --- Summary Table ---
\subsubsection{Summary: Neural Architectures for Sequences}

\begin{center}
\renewcommand{\arraystretch}{1.3}
\begin{tabular}{@{} l p{3cm} p{3cm} p{4cm} @{}}
\toprule
\textbf{Architecture} & \textbf{Context} & \textbf{Key Feature} & \textbf{Use Cases} \\
\midrule
Feedforward NN & Fixed window & Simple, parallelisable & Classification, fixed-context LM \\
Simple RNN & Unbounded (via $h_{t-1}$) & Sequential processing & LM, sequence tasks \\
LSTM & Unbounded + gated & Handles long-range dependencies & LM, translation, tagging \\
Encoder-Decoder & Full source sequence & Seq-to-seq mapping & Translation, summarisation \\
+ Attention & Dynamic focus & Soft alignment & Improved translation \\
\bottomrule
\end{tabular}
\end{center}

